% !TEX program = xelatex
\documentclass[11pt,a4paper]{article}

\usepackage[UTF8]{ctex}
\usepackage[margin=2.2cm]{geometry}
\usepackage{amsmath, amssymb, amsthm}
\usepackage{listings}
\usepackage{xcolor}
\usepackage{tikz}
\usepackage{booktabs}
\usepackage{multirow}
\usepackage{longtable}
\usepackage{seqsplit}
\usepackage{hyperref}
\usetikzlibrary{positioning, arrows.meta, shapes.geometric, fit, calc, decorations.pathreplacing}

\hypersetup{
  colorlinks=true,
  linkcolor=blue!60!black,
  urlcolor=teal!70!black,
}

\lstdefinelanguage{Rust}{
  morekeywords={fn,let,mut,pub,struct,impl,use,mod,self,Self,Option,Some,None,return,if,else,for,while,loop,match,break,continue,as,where,trait,type,enum,move,ref,const,static,unsafe,extern,crate,super,dyn,async,await,box,in},
  sensitive=true,
  morecomment=[l]{//},
  morecomment=[s]{/*}{*/},
  morestring=[b]",
}

\lstset{
  language=Rust,
  basicstyle=\ttfamily\footnotesize,
  keywordstyle=\color{blue!60!black}\bfseries,
  commentstyle=\color{green!40!black}\itshape,
  stringstyle=\color{red!60!black},
  numbers=left,
  numberstyle=\tiny\color{gray},
  numbersep=8pt,
  frame=single,
  framerule=0.4pt,
  rulecolor=\color{gray!40},
  breaklines=true,
  breakatwhitespace=true,
  showstringspaces=false,
  tabsize=4,
  columns=fullflexible,
  keepspaces=true,
}

\title{\textbf{halfedge\,: 测地线距离模块设计文档} \\ \large \texttt{src/geodesics.rs}}
\author{halfedge 库}
\date{2026-07-05}

\begin{document}
\maketitle
\tableofcontents

\section{模块定位}

\texttt{geodesics.rs} 在 \texttt{geometry.rs} / \texttt{linalg.rs} 之上实现三种测地线距离算法：
\begin{itemize}
  \item \textbf{Heat Method}（Crane, Weischedel, Wardetzky 2013）— 近似测地线
  \item \textbf{Dijkstra} — 图上最短路径（严格上界）
  \item \textbf{MMP}（Mitchell-Mount-Papadimitriou 1987）— 精确测地线
\end{itemize}
并提供沿距离场梯度回溯最短路径的能力。

\begin{center}
\renewcommand{\arraystretch}{1.18}
\begin{tabular}{@{}llp{6.0cm}@{}}
  \toprule
  \textbf{类别} & \textbf{API} & \textbf{语义} \\
  \midrule
  距离场（Heat Method）
    & \texttt{geodesic\_distance\_from\_vertex}
    & Heat Method：从 \texttt{source} 出发计算所有顶点的测地线距离；
      返回顺序对应 \texttt{mesh.vertex\_ids()}；求解失败返回 \texttt{None}。 \\
  路径回溯
    & \texttt{shortest\_path}
    & 沿距离场梯度回溯最短路径：返回从 \texttt{target} 到源顶点的顶点序列（含两端）。 \\
  \midrule
  距离场（Dijkstra）
    & \texttt{dijkstra\_geodesic}
    & 单源 Dijkstra 最短\textbf{图}距离（边长权值），不依赖线性求解器 \\
  距离场（Dijkstra 多源）
    & \texttt{dijkstra\_multi\_source\_geodesic}
    & 多源 Dijkstra：每个顶点到最近源点的图距离 \\
  路径回溯（Dijkstra）
    & \texttt{dijkstra\_with\_parent}
    & Dijkstra + 父节点回溯数组，用于精确路径重建 \\
  路径（Dijkstra 精确）
    & \texttt{dijkstra\_shortest\_path}
    & 返回 \texttt{source} 到 \texttt{target} 的最短顶点序列（含两端） \\
  \midrule
  距离场（多源 Heat Method）
    & \texttt{multi\_source\_geodesic}
    & 多源 Heat Method：每个顶点到最近源点的测地线距离；
      源点强制为 $0$ 消除数值误差 \\
  \midrule
  距离场（MMP 精确）
    & \texttt{mmp\_geodesic}
    & MMP 精确测地线：从 \texttt{source} 出发计算所有顶点的精确测地线距离；
      可穿过三角面内部，非仅沿边 \\
  距离场（MMP 多源）
    & \texttt{mmp\_multi\_source\_geodesic}
    & 多源 MMP：每个顶点到最近源点的精确测地线距离 \\
  \bottomrule
\end{tabular}
\end{center}

\section{Heat Method 算法}

Heat Method 把\textbf{非线性}的测地线距离计算转化为\textbf{两个线性系统}：
先用热流方程扩散初始脉冲得到光滑梯度方向，再用 Poisson 方程从散度场恢复距离。
其核心洞察是：\emph{热流梯度的方向与测地线梯度方向几乎一致}，且方向对参数不敏感。

\subsection{Step 1：时间步 \texorpdfstring{$\tau = h^2$}{tau = h\^{}2}}

设网格平均边长为 $h$，时间步取
\[
  \tau = h^2, \qquad
  h = \frac{1}{|E|}\sum_{e \in E} L(e).
\]
\textbf{取 $h^2$ 的理由}：使热流在每条边上的扩散量
$\tau \cdot \frac{1}{h^2} \sim O(1)$，既不会过快抹平局部信息，也不会过慢导致数值噪声。
代码中若 $|E| = 0$，回退 $\tau = 1$。

\subsection{Step 2：热流扩散}

求解线性系统
\[
  (M + \tau L)\, u_\tau = M\, u_0,
\]
其中 $M$ 为 lumped mass（顶点 Voronoi 面积对角矩阵），$L$ 为余切拉普拉斯，
$u_0$ 为初始热源。初始热源取源顶点 $s$ 处的\textbf{单位积分脉冲}：
\[
  u_0[i] =
  \begin{cases}
    1 / M_{ss}, & i = s, \\
    0,          & \text{otherwise}.
  \end{cases}
\]
等价于热传导方程 $\frac{\partial u}{\partial t} = \Delta u$ 在时间 $\tau$ 后的解。
代码使用\textbf{共轭梯度法}（\texttt{conjugate\_gradient}）求解，
最大迭代数 $100\,n$，相对残差 $10^{-6}$，并预先调用
\texttt{regularize\_diagonal} 加入 $10^{-10}$ 对角正则化以保证 SPD 性。

\subsection{Step 3：单位负梯度}

对每个三角面 $f$，由分段线性基函数求 $u_\tau$ 的常梯度
$\nabla u_\tau|_f$，再归一化并取负：
\[
  X_f = -\frac{\nabla u_\tau|_f}{\|\nabla u_\tau|_f\|}.
\]
\textbf{取负原因}：热流从高温（源）流向低温（远端），
$\nabla u_\tau$ 指向\emph{远离源}，而测地线最短路径应\emph{朝向源}前进，故取负。
退化面（$\|\nabla u_\tau\| < 10^{-10}$）置 $X_f = \vec{0}$，避免除零。

\subsection{Step 4：散度与 Poisson 求解}

将单位负梯度场 $X$ 投影回顶点散度（弱形式）：
\[
  (\nabla \cdot X)_i = \sum_{f \ni i} A_f \cdot \bigl( X_f \cdot \nabla \varphi_i|_f \bigr),
\]
其中 $\varphi_i$ 是顶点 $i$ 的分片线性基函数。
再求解 Poisson 方程
\[
  L\, \phi = \nabla \cdot X,
\]
$L$ 与 Step 2 共用\textbf{同一份}余切拉普拉斯矩阵 \texttt{cot\_lap}（\texttt{CsMat<f64>}）：
Step 1 通过 \texttt{build\_laplacian\_and\_mass} 构建一次，Step 2 用其缩放 $\tau L$
加入 mass 对角得 $M + \tau L$，Step 4 直接 \texttt{clone} 后做对角正则化用于 Poisson 求解。
\textbf{避免重复构建}拉普拉斯（旧实现构建 3 次：Step 2、散度函数内部、Step 4 各一次）。
同样使用共轭梯度法求解。

\subsection{Step 5：距离场偏移校正}

Poisson 解 $\phi$ 与真实测地线距离 $d$ 仅差一个常数（$\Delta$ 的核空间）。
取源顶点 $s$ 处为零参考点：
\[
  d_i = \bigl|\, \phi_i - \phi_s \,\bigr|.
\]
取绝对值是防止数值误差导致个别顶点 $\phi_i$ 略小于 $\phi_s$。
校正后 $d_s = 0$，其余 $d_i \ge 0$。

\section{Heat Method 算法流程图}

\begin{center}
\resizebox{\textwidth}{!}{%
\begin{tikzpicture}[
  node distance=0.55cm,
  proc/.style={draw, rounded corners=2pt, minimum width=10.5cm, minimum height=0.75cm, align=center, font=\small},
  startend/.style={draw, rounded corners=10pt, minimum width=2.8cm, minimum height=0.7cm, font=\small, fill=gray!12},
  op/.style={-Stealth, thick},
]
  \node[startend] (s) {开始：\texttt{source} 已知};
  \node[proc, below=of s, fill=blue!8] (p1)
    {Step 1 计算 $h = \frac{1}{|E|}\sum L(e)$，$\tau = h^2$};
  \node[proc, below=of p1, fill=blue!12] (p2)
    {Step 2 构建并求解 $(M + \tau L)\, u_\tau = M\, u_0$（CG，残差 $10^{-6}$）};
  \node[proc, below=of p2, fill=green!12] (p3)
    {Step 3 逐面求 $\nabla u_\tau|_f$ 并归一化 $X_f = -\nabla u_\tau / \|\nabla u_\tau\|$};
  \node[proc, below=of p3, fill=yellow!18] (p4)
    {Step 4 弱形式积分 $\nabla \cdot X$，求解 $L\, \phi = \nabla \cdot X$（CG）};
  \node[proc, below=of p4, fill=red!10] (p5)
    {Step 5 偏移校正 $d_i = |\phi_i - \phi_s|$};
  \node[startend, below=of p5] (e) {返回 \texttt{Some}($d$)};
  \draw[op] (s) -- (p1);
  \draw[op] (p1) -- (p2);
  \draw[op] (p2) -- (p3);
  \draw[op] (p3) -- (p4);
  \draw[op] (p4) -- (p5);
  \draw[op] (p5) -- (e);
\end{tikzpicture}
}
\end{center}

\section{离散算子}

\subsection{离散拉普拉斯（余切权重）}

对相邻顶点 $i, j$，设边 $ij$ 的两个对角为 $\alpha_{ij}, \beta_{ij}$，则
\[
  L_{ij} = -\frac{1}{2}\bigl(\cot \alpha_{ij} + \cot \beta_{ij}\bigr),
  \qquad
  L_{ii} = -\sum_{j \ne i} L_{ij}.
\]
\textbf{边界边}只有一个对角，使用单边余切。
实现中仅累加 $w > 0$ 的项（非钝角三角形贡献）以避免负权重导致矩阵非 SPD。
Lumped mass $M_{ii}$ 取顶点周围面面积的 $1/3$ 之和，退化时钳为 $10^{-14}$。

\subsection{面梯度（FEM 形函数）}

对三角面 $f$ 的三顶点 $i, j, k$，分片线性函数
$u|_f = u_i \varphi_i + u_j \varphi_j + u_k \varphi_k$ 的常梯度为
\[
  \nabla u|_f = \frac{1}{2 A_f}\, \hat{n}_f \times
  \bigl( u_i \vec{e}_i + u_j \vec{e}_j + u_k \vec{e}_k \bigr),
\]
其中 $\vec{e}_i = \vec{p}_k - \vec{p}_j$ 是顶点 $i$ 所对的边向量，
$A_f$ 是面面积，$\hat{n}_f$ 是面法向。

\subsection{面散度（弱形式积分）}

向量场 $X$ 在顶点 $i$ 处的散度（弱形式）为
\[
  (\nabla \cdot X)_i = \sum_{f \ni i} A_f \cdot \bigl( X_f \cdot \nabla \varphi_i|_f \bigr),
  \qquad
  \nabla \varphi_i|_f = \frac{1}{2 A_f}\, \hat{n}_f \times \vec{e}_i,
\]
等价地写作
\[
  (\nabla \cdot X)_i = \frac{1}{2} \sum_{f \ni i} X_f \cdot \bigl( \hat{n}_f \times \vec{e}_i \bigr).
\]
代码中 $\nabla \varphi_i|_f$ 直接用 $\hat{n}_f \times \vec{e}_i / (2 A_f)$ 计算，
再乘以 $A_f$ 与 $X_f$ 的点积，与上式一致。

\section{最短路径回溯}

\texttt{shortest\_path} 在已计算好的距离场 $d$ 上，沿\textbf{贪心梯度下降}
方向回溯：当前顶点 $v$ 的邻域中，选择 $d$ 最小且严格小于当前距离的顶点 $v'$，
跳转至 $v'$，重复直至到达源点（$d < 10^{-10}$）或陷入局部极小（无更小邻点）。
步数上界 $2|V|$ 防止死循环。

\begin{center}
\begin{tikzpicture}[
  >=Stealth,
  vert/.style={circle, draw=blue!70!black, fill=blue!10, minimum size=8mm, inner sep=0pt, font=\footnotesize},
  src/.style={circle, draw=red!70!black, fill=red!12, minimum size=8mm, inner sep=0pt, font=\footnotesize},
  arr/.style={->, thick, red!70!black},
]
  % 顶点
  \node[src] (s) at (0,0) {$s$};
  \node[vert] (a) at (1.6,1.0) {$a$};
  \node[vert] (b) at (1.6,-1.0) {$b$};
  \node[vert] (c) at (3.4,1.6) {$c$};
  \node[vert] (d) at (3.4,-1.6) {$d$};
  \node[vert] (e) at (5.0,0.6) {$e$};
  \node[vert] (t) at (6.6,0) {$t$};
  % 边
  \draw[gray!50] (s) -- (a);
  \draw[gray!50] (s) -- (b);
  \draw[gray!50] (a) -- (c);
  \draw[gray!50] (a) -- (e);
  \draw[gray!50] (b) -- (d);
  \draw[gray!50] (b) -- (e);
  \draw[gray!50] (c) -- (e);
  \draw[gray!50] (d) -- (e);
  \draw[gray!50] (e) -- (t);
  \draw[gray!50] (c) -- (t);
  % 回溯路径
  \draw[arr] (t) -- (e);
  \draw[arr] (e) -- (a);
  \draw[arr] (a) -- (s);
  % 距离标签
  \node[font=\tiny, gray] at (-0.5,-0.5) {$d=0$};
  \node[font=\tiny, gray] at (1.6,1.5) {$d=1.1$};
  \node[font=\tiny, gray] at (5.0,1.1) {$d=2.4$};
  \node[font=\tiny, gray] at (6.6,0.5) {$d=3.3$};
  % 标题
  \node[font=\footnotesize, align=center] at (3,-2.4)
    {回溯方向：$t \to e \to a \to s$（沿 $d$ 严格递减）};
\end{tikzpicture}
\end{center}

\textbf{贪心策略的局限}：若距离场存在局部极小（数值误差或网格各向异性），
可能在到达源点前提前终止；上界 $2|V|$ 保证不进入死循环。
返回序列包含两端顶点，从 \texttt{target} 起到源点止。

\section{Dijkstra 最短图距离}

\texttt{dijkstra\_geodesic} 使用经典 Dijkstra 算法计算\textbf{图最短距离}：
以边长（欧氏长度）为权值，求源点到所有顶点的最短路径长度。
相比 Heat Method，Dijkstra 路径只能沿网格边行进，结果为\textbf{图距离上界}，
但实现简单、不依赖线性求解器，且结果\textbf{精确}（无数值误差）。

\subsection{最小堆实现}

Rust 标准库 \texttt{BinaryHeap} 是最大堆，通过反转 \texttt{Ord} 比较方向实现最小堆：
\begin{lstlisting}
#[derive(Clone, Copy, Debug, PartialEq)]
struct QueueEntry { dist: f64, vertex: usize }

impl Ord for QueueEntry {
    fn cmp(&self, other: &Self) -> Ordering {
        // 比较方向取反，使 BinaryHeap 成为最小堆
        other.dist.partial_cmp(&self.dist).unwrap_or(Ordering::Equal)
            .then_with(|| other.vertex.cmp(&self.vertex))
    }
}
\end{lstlisting}

\subsection{算法流程}

\begin{enumerate}
  \item 初始化 \texttt{dist[source] = 0}，其余为 $+\infty$；
  \item 将 $(0, \text{source})$ 入堆；
  \item 循环弹出堆顶 $(d, v)$：
    \begin{itemize}
      \item 若 $d > \text{dist}[v]$，跳过（已通过更短路径访问过）；
      \item 遍历 $v$ 的邻居 $u$，若 $d + \|p_u - p_v\| < \text{dist}[u]$，
            更新 $\text{dist}[u]$ 并入堆 $(d + \|p_u - p_v\|, u)$；
    \end{itemize}
  \item 堆空时终止。
\end{enumerate}

复杂度 $O(|E|\log|V|)$。

\paragraph{顶点索引反查的常数优化}
主循环中需要把堆里存的索引 $u$ 反查回 \texttt{VertexId}。
早期实现调用 \texttt{mesh.vertex\_ids().nth(u).unwrap()}：\texttt{nth} 每次从
迭代器头部推进 $u$ 步，单次 $O(u)$，整体退化为 $O(|V|^2)$。
修复方案：入口处一次性构建 \texttt{vid\_list: Vec<VertexId>}（与
\texttt{v\_idx: HashMap<VertexId, usize>} 同步生成，单次遍历），
循环内用 \texttt{vid\_list[u]} 做 $O(1)$ 数组下标访问。
\texttt{dijkstra\_geodesic}、\texttt{dijkstra\_multi\_source\_geodesic}、
\texttt{dijkstra\_with\_parent} 三个入口均已采用此模式。

\subsection{多源 Dijkstra}

\texttt{dijkstra\_multi\_source\_geodesic} 将所有源点以距离 $0$ 同时入堆，
自然得到每个顶点到\textbf{最近源点}的距离。

\subsection{路径回溯：dijkstra\_with\_parent}

\texttt{dijkstra\_with\_parent} 在 Dijkstra 同时维护 \texttt{parent} 数组：
每次松弛 \texttt{dist[u]} 时记录 \texttt{parent[u] = v}。
返回 \texttt{(dist, parent)} 二元组，便于精确路径重建。

\texttt{dijkstra\_shortest\_path} 在此基础上从 \texttt{target} 沿
\texttt{parent} 链回溯到 \texttt{source}，得到精确最短顶点序列。

\begin{center}
\begin{tikzpicture}[
  >=Stealth,
  node distance=0.4cm,
  proc/.style={draw, rounded corners=2pt, minimum width=10cm, minimum height=0.65cm, align=center, font=\small},
  startend/.style={draw, rounded corners=10pt, minimum width=2.6cm, minimum height=0.65cm, font=\small, fill=gray!12},
  op/.style={-Stealth, thick}
]
  \node[startend] (s) {开始：\texttt{source}, \texttt{target}};
  \node[proc, below=of s, fill=blue!8] (p1)
    {运行 \texttt{dijkstra\_with\_parent}，得到 \texttt{dist} 与 \texttt{parent}};
  \node[proc, below=of p1, fill=yellow!18] (p2)
    {\texttt{cur = target}，\texttt{path = [target]}};
  \node[proc, below=of p2, fill=green!12] (p3)
    {循环：\texttt{path.push(parent[cur])}，\texttt{cur = parent[cur]}，
     直到 \texttt{cur == source} 或 \texttt{parent[cur]} 无定义};
  \node[proc, below=of p3, fill=orange!12] (p4)
    {\texttt{path.reverse()}（使顺序为 \texttt{source} $\to$ \texttt{target}）};
  \node[startend, below=of p4] (e) {返回 \texttt{path}};
  \draw[op] (s) -- (p1);
  \draw[op] (p1) -- (p2);
  \draw[op] (p2) -- (p3);
  \draw[op] (p3) -- (p4);
  \draw[op] (p4) -- (e);
\end{tikzpicture}
\end{center}

\section{多源 Heat Method}

\texttt{multi\_source\_geodesic} 将 Heat Method 扩展到\textbf{多源}场景：
计算每个顶点到最近源点的测地线距离。实现上仅修改初始热源 $u_0$：
\[
  u_0[i] =
  \begin{cases}
    1 / M_{ii}, & i \in S \text{（任一源点）}, \\
    0,          & \text{otherwise}.
  \end{cases}
\]
其余步骤（热流扩散、梯度归一化、Poisson 求解、偏移校正）与单源完全相同，
且\textbf{共用同一份} \texttt{cot\_lap}：构建一次后热流系统与 Poisson 系统分别引用
（Poisson 求 \texttt{clone} 后做正则化），避免重复构建。

\textbf{源点数值误差处理}：Heat Method 在源点处存在数值误差，
单源时通过 $\phi_s$ 偏移校正消除；多源时各源点误差不一致，
故在偏移校正后\textbf{显式将每个源点距离置零}：
\[
  \forall s \in S: \quad d_s \leftarrow 0.
\]

\section{MMP 精确测地线算法}

\texttt{mmp\_geodesic} 实现 Mitchell-Mount-Papadimitriou (1987) 精确测地线算法，
通过\textbf{窗口传播}（window propagation）在三角面上追踪伪波前，
计算从源顶点到所有顶点的精确测地线距离。
与 Dijkstra 不同，MMP 的测地线路径可\textbf{穿过三角面内部}，
而非仅沿边行进；与 Heat Method 不同，MMP 给出\textbf{精确}距离而非近似值。

\subsection{核心概念：窗口（Window）}

窗口是三角面边上的一个\textbf{伪波前段}，记录了波前到达边上一个区间的信息：

\begin{center}
\renewcommand{\arraystretch}{1.18}
\begin{tabular}{@{}lp{9.0cm}@{}}
  \toprule
  \textbf{字段} & \textbf{语义} \\
  \midrule
  \texttt{he} & 窗口所在的半边 \\
  \texttt{b0}, \texttt{b1} & 窗口在边上的参数范围 $[b_0, b_1] \subset [0, 1]$；
    0 = origin（\texttt{he.twin.vertex}），1 = tip（\texttt{he.vertex}） \\
  \texttt{d0}, \texttt{d1} & $b_0$ 和 $b_1$ 处的测地线距离 \\
  \texttt{pseudo\_src} & 伪源位置（3D），经展开反射后；
    窗口内任意点 $x$ 的距离为 $\|S' - x\|$ \\
  \texttt{from\_face} & 窗口来自的面（传播方向：进入另一个面） \\
  \bottomrule
\end{tabular}
\end{center}

窗口内任意点 $x$ 的测地线距离等于伪源到该点的欧氏距离：
\[
  d(x) = \|S' - x\|.
\]
这一性质是 MMP 算法正确性的基础：展开操作保证从原始源到 $x$ 的
测地线路径在展开后成为直线段。

\subsection{展开（Unfolding）}

当窗口跨过共享边从源面传播到目标面时，需要将伪源从源面平面
``展开''到目标面平面。展开操作保持伪源到共享边上各点距离不变，
同时使伪源到目标面内各点的 3D 距离等于测地线距离。

\begin{center}
\begin{tikzpicture}[
  >=Stealth,
  node distance=0.5cm,
  face/.style={draw, thick, fill=#1, minimum width=3cm, minimum height=2cm, align=center, font=\small},
  arr/.style={-Stealth, thick, red!70!black},
  lbl/.style={font=\footnotesize},
]
  % 源面
  \node[face=blue!8] (f1) at (0,0) {源面 $f_1$};
  \node[lbl, above=0.1cm of f1] {$S$（伪源）};
  % 共享边
  \node[draw, thick, fill=orange!20, minimum width=0.3cm, minimum height=2cm] (edge) at (3.2,0) {};
  \node[lbl, below=0.1cm of edge] {共享边 $e$};
  % 目标面
  \node[face=green!8] (f2) at (6.4,0) {目标面 $f_2$};
  \node[lbl, above=0.1cm of f2] {$S'$（展开后）};
  % 箭头
  \draw[arr] (3.5,0.5) -- (5.0,0.5) node[midway, above, font=\footnotesize] {展开};
  % 旋转示意
  \draw[dashed, gray] (3.2,1.3) arc (90:70:1.3) node[midway, above, font=\tiny, gray] {旋转};
\end{tikzpicture}
\end{center}

实现中，展开通过构建局部坐标系完成：
\begin{enumerate}
  \item 以共享边为 $x$ 轴构建局部坐标系；
  \item 在源面坐标系中表达伪源坐标 $(s_x, s_y)$；
  \item 将坐标映射到目标面坐标系（更换 $y$ 轴方向为目标面法向侧）。
\end{enumerate}

\subsection{窗口传播}

给定窗口 $w$ 位于边 $e$ 上，\texttt{from\_face} $= f_1$，传播步骤如下：

\begin{enumerate}
  \item \textbf{确定目标面}：$f_2$ 为边 $e$ 的另一个邻接面（$\ne f_1$）；
  \item \textbf{获取对面顶点}：$f_2$ 的三个顶点中不在 $e$ 上的那个顶点 $C$；
  \item \textbf{展开伪源}：$S' = \text{unfold}(S, f_1, f_2, e)$；
  \item \textbf{更新顶点距离}：$d_C = \min(d_C, \|S' - C\|)$；
  \item \textbf{计算子窗口}：伪源 $S'$ 通过窗口 $[b_0, b_1]$ 发出的射线楔
    与两条子边（$O \to C$ 和 $T \to C$）的交集。
\end{enumerate}

\subsection{子窗口计算}

子窗口的计算在 2D 局部坐标系中进行（投影到面平面后保持等距性）：

\begin{center}
\begin{tikzpicture}[
  >=Stealth,
  scale=1.2,
  vert/.style={circle, draw=blue!70!black, fill=blue!10, minimum size=6mm, inner sep=0pt, font=\footnotesize},
  src/.style={circle, draw=red!70!black, fill=red!12, minimum size=6mm, inner sep=0pt, font=\footnotesize},
  win/.style={line width=2.5pt, red!60!black},
  newwin/.style={line width=2.5pt, green!50!black},
  wedge/.style={->, dashed, orange!70!black, thick},
]
  % 三角形顶点
  \node[vert] (O) at (0,0) {$O$};
  \node[vert] (T) at (4,0) {$T$};
  \node[vert] (C) at (2,3) {$C$};
  % 伪源
  \node[src] (S) at (2,-1.5) {$S'$};
  % 窗口
  \draw[win] (1.2,0) -- (2.8,0) node[midway, below, font=\footnotesize] {窗口 $[b_0, b_1]$};
  % 楔射线
  \draw[wedge] (S) -- (1.2,0);
  \draw[wedge] (S) -- (2.8,0);
  % 子窗口
  \draw[newwin] (0.3,0.45) -- (1.5,2.25) node[midway, left, font=\footnotesize] {子窗口$_1$};
  \draw[newwin] (2.8,0) -- (1.8,1.5) node[midway, right, font=\footnotesize] {子窗口$_2$};
  % 三角形边
  \draw[gray!50] (O) -- (T);
  \draw[gray!50] (T) -- (C);
  \draw[gray!50] (C) -- (O);
  % 标注
  \draw[decorate, decoration={brace, amplitude=5pt, mirror}] (0,-0.3) -- (4,-0.3)
    node[midway, below=6pt, font=\footnotesize] {边 $e$（源面 $\to$ 目标面）};
\end{tikzpicture}
\end{center}

具体步骤：
\begin{enumerate}
  \item 将伪源 $S'$、窗口端点 $P(b_0)$/$P(b_1)$、子边端点投影到面平面 2D 坐标系；
  \item 构建照明楔：从 $S'$ 出发通过 $P(b_0)$ 和 $P(b_1)$ 的两条射线；
  \item 对每条子边，计算楔与子边的交集：
    \begin{itemize}
      \item 检查子边端点是否在楔内（2D 叉积判定）；
      \item 计算楔边界射线与子边的交点（2D 射线-线段求交）；
      \item 合并所有交点得到参数区间 $[s_{\min}, s_{\max}]$；
    \end{itemize}
  \item 在 3D 中计算子窗口端点的距离：$d = \|S' - Q\|$（使用展开后的伪源）；
  \item 创建新的 \texttt{MmpWindow} 加入优先队列。
\end{enumerate}

\subsection{算法流程图}

\begin{center}
\resizebox{\textwidth}{!}{%
\begin{tikzpicture}[
  node distance=0.5cm,
  proc/.style={draw, rounded corners=2pt, minimum width=11cm, minimum height=0.75cm, align=center, font=\small},
  startend/.style={draw, rounded corners=10pt, minimum width=2.8cm, minimum height=0.7cm, font=\small, fill=gray!12},
  decision/.style={draw, diamond, aspect=2.5, minimum width=3cm, align=center, font=\small, fill=yellow!15},
  op/.style={-Stealth, thick},
]
  \node[startend] (s) {开始：\texttt{source} 已知};
  \node[proc, below=of s, fill=blue!8] (p1)
    {初始化：对 source 的每个邻接面，在对面边创建窗口 $[0,1]$，伪源 $= \text{source}$};
  \node[proc, below=of p1, fill=blue!12] (p2)
    {将所有初始窗口按 $\min(d_0, d_1)$ 加入最小堆};
  \node[decision, below=of p2] (d1) {堆空？};
  \node[proc, below=of d1, fill=orange!12] (p3)
    {弹出最小距离窗口 $w$；剪枝：若 $\min(d_0,d_1)$ 大于边端点已知距离 $+ \varepsilon$，跳过};
  \node[proc, below=of p3, fill=green!12] (p4)
    {展开伪源 $S'$：从源面坐标系映射到目标面坐标系};
  \node[proc, below=of p4, fill=green!8] (p5)
    {更新对面顶点距离：$d_C = \min(d_C, \|S' - C\|)$};
  \node[proc, below=of p5, fill=yellow!18] (p6)
    {2D 投影 + 楔交集计算子窗口 $[s_{\min}, s_{\max}]$（两条子边）};
  \node[proc, below=of p6, fill=red!10] (p7)
    {3D 计算子窗口距离，创建新 \texttt{MmpWindow} 入堆};
  \node[decision, below=of p7] (d2) {迭代 $> |F| \times 100$？};
  \node[startend, below=of d2] (e) {返回 $d$};

  \draw[op] (s) -- (p1);
  \draw[op] (p1) -- (p2);
  \draw[op] (p2) -- (d1);
  \draw[op] (d1) -- node[left, font=\footnotesize] {否} (p3);
  \draw[op] (p3) -- (p4);
  \draw[op] (p4) -- (p5);
  \draw[op] (p5) -- (p6);
  \draw[op] (p6) -- (p7);
  \draw[op] (p7) -- (d2);
  \draw[op] (d2) -- node[left, font=\footnotesize] {否} ++(-3,0) |- (d1);
  \draw[op] (d1) -- node[above, font=\footnotesize] {是} ++(5,0) |- (e);
  \draw[op] (d2) -- node[above, font=\footnotesize] {是} (e);
\end{tikzpicture}
}
\end{center}

\subsection{多源 MMP：统一传播}

\texttt{mmp\_multi\_source\_geodesic} 与单源 \texttt{mmp\_geodesic} 共用同一核心
\texttt{mmp\_geodesic\_impl}。多源扩展的关键是\textbf{初始化阶段}：

\begin{enumerate}
  \item 遍历所有源 $s \in S$，对每个 $s$ 的邻接面创建初始窗口；
  \item 所有窗口\textbf{放入同一个优先队列}；
  \item 主循环与单源\textbf{完全相同}——窗口传播自然处理来自不同源的波前相遇。
\end{enumerate}

\paragraph{为何不逐源运行取最小？}
旧实现「\texttt{for s in S: d\_s = mmp(s); result = min(result, d\_s)}」复杂度为
$O(S \cdot n^2 \log n)$，随源数 $S$ 线性增长。统一传播复杂度仍为
$O(n^2 \log n)$（窗口总数随 $S$ 略增但同阶），\textbf{不}随 $S$ 线性退化。

\paragraph{正确性}
统一传播得到的距离场 $d_{\text{unified}}[v]$ 与逐源取最小
$d_{\text{naive}}[v] = \min_{s \in S} d_s[v]$ 在浮点容差内\textbf{数值一致}
（测试 \texttt{mmp\_multi\_source\_unified\_matches\_naive} 验证）。
原因：MMP 窗口的伪源位置 \texttt{pseudo\_src} 已编码源信息，
不同源的窗口在传播时互不干扰，相遇时距离较小者自然胜出。

\subsection{复杂度}

最坏情况 $O(n^2 \log n)$，实践中通常远好于此。
剪枝策略（跳过最小距离大于已知顶点距离的窗口）有效减少了实际处理的窗口数量。

\section{三种算法对比}

\begin{center}
\renewcommand{\arraystretch}{1.18}
\begin{tabular}{@{}p{2.8cm}p{3.8cm}p{3.8cm}p{3.8cm}@{}}
  \toprule
  \textbf{属性} & \textbf{Heat Method} & \textbf{Dijkstra} & \textbf{MMP} \\
  \midrule
  路径类型 & 测地线（可跨面） & 图路径（仅沿边） & 测地线（可跨面） \\
  精度 & 近似（数值求解） & 图精确（无误差） & 精确（无离散化误差） \\
  依赖 & 稀疏线性求解器（CG） & 仅需最小堆 & 仅需最小堆 \\
  复杂度 & $O(k \cdot n)$（CG 迭代） & $O(|E|\log|V|)$ & $O(n^2 \log n)$ 最坏 \\
  源点误差 & 有（偏移校正后 $\approx 0$） & 无 & 无 \\
  多源扩展 & 修改 $u_0$ + 显式置零 & 多源同时入堆 & 多源同时入堆（统一传播） \\
  路径提取 & 梯度回溯（近似） & parent 回溯（精确图路径） & 梯度回溯（近似） \\
  \bottomrule
\end{tabular}
\end{center}

\textbf{关系}：
\begin{itemize}
  \item Dijkstra 距离是真实测地线距离的\textbf{上界}（仅沿边，路径更长）；
  \item MMP 距离是真实测地线距离的\textbf{精确值}（可穿面）；
  \item Heat Method 距离是真实测地线距离的\textbf{近似值}（数值求解）；
  \item 因此 $d_{\text{MMP}} \le d_{\text{Dijkstra}}$ 恒成立。
\end{itemize}

\section{退化守卫}

\begin{center}
\renewcommand{\arraystretch}{1.18}
\begin{tabular}{@{}lp{4.0cm}p{5.4cm}@{}}
  \toprule
  \textbf{函数} & \textbf{退化场景} & \textbf{处理} \\
  \midrule
  \texttt{geodesic\_distance\_from\_vertex}
    & 空网格（$|V| = 0$） & 返回 \texttt{Some}(\texttt{Vec::new()}) \\
  \texttt{geodesic\_distance\_from\_vertex}
    & \texttt{source} 不在 \texttt{v\_idx} 中 & 返回 \texttt{None}（\texttt{?} 传播） \\
  \texttt{geodesic\_distance\_from\_vertex}
    & $|E| = 0$ & 时间步回退 $\tau = 1$ \\
  \texttt{geodesic\_distance\_from\_vertex}
    & 共轭梯度未收敛 & CG 返回 \texttt{None}，函数返回 \texttt{None} \\
  \texttt{build\_laplacian\_and\_mass}
    & 顶点面积 $< 10^{-14}$ & 钳为 $10^{-14}$，防除零 \\
  \texttt{build\_laplacian\_and\_mass}
    & 余切权重 $w \le 0$ & 跳过该边，避免破坏 SPD \\
  \texttt{compute\_face\_gradients}
    & 面积 $< 10^{-14}$（退化面） & 跳过该面 \\
  Step 3 归一化
    & $\|\nabla u_\tau\| < 10^{-10}$ & $X_f = \vec{0}$ \\
  \texttt{shortest\_path}
    & \texttt{target} 无效或越界 & 返回 \texttt{Vec::new()} \\
  \texttt{shortest\_path}
    & 陷入局部极小 & 提前终止，返回已收集路径 \\
  \texttt{shortest\_path}
    & 步数达上界 $2|V|$ & 终止循环，防止死循环 \\
  \midrule
  \texttt{mmp\_geodesic}
    & 空网格 & 返回空 \texttt{Vec} \\
  \texttt{mmp\_geodesic}
    & \texttt{source} 无效 & 返回全 \texttt{INFINITY} \\
  \texttt{mmp\_geodesic}
    & 面非三角（$|he| \ne 3$） & 跳过该面 \\
  \texttt{mmp\_geodesic}
    & 迭代超限（$> |F| \times 100$） & 终止循环 \\
  \texttt{unfold\_pseudo\_source}
    & 退化边（长度 $< 10^{-14}$） & 返回原始伪源（不反射） \\
  \texttt{propagate\_window}
    & 边界半边（无 twin） & 返回空窗口列表 \\
  \texttt{propagate\_window}
    & 子边长度 $< 10^{-14}$ & 跳过该子边 \\
  \texttt{propagate\_window}
    & 窗口区间 $< 10^{-10}$ & 跳过（退化窗口） \\
  \bottomrule
\end{tabular}
\end{center}

\section{使用示例}

\begin{lstlisting}
use halfedge::geodesics::{
    geodesic_distance_from_vertex, shortest_path,
    dijkstra_geodesic, mmp_geodesic,
};
use halfedge::test_util::build_icosphere;

let mesh = build_icosphere(2);
let vertices: Vec<_> = mesh.vertex_ids().collect();
let source = vertices[0];
let target = vertices[vertices.len() / 2];

// 1. Heat Method（近似）
let dist_heat = geodesic_distance_from_vertex(&mesh, source)
    .expect("heat method ok");

// 2. Dijkstra（图距离，上界）
let dist_dijk = dijkstra_geodesic(&mesh, source);

// 3. MMP（精确测地线）
let dist_mmp = mmp_geodesic(&mesh, source);

// 精确距离 <= 图距离
assert!(dist_mmp[target_idx] <= dist_dijk[target_idx] + 1e-8);

// 4. 从 target 回溯最短路径
let path = shortest_path(&mesh, &dist_mmp, target);
\end{lstlisting}

\section{测试覆盖}

{\footnotesize
\renewcommand{\arraystretch}{1.18}
\renewcommand{\_}{\textunderscore\allowbreak}%
\begin{longtable}{@{}p{5.6cm}p{9.6cm}@{}}
  \toprule
  \textbf{测试名} & \textbf{验证点} \\
  \midrule
  \endfirsthead
  \multicolumn{2}{c}{\small （续表）} \\
  \toprule
  \textbf{测试名} & \textbf{验证点} \\
  \midrule
  \endhead
  \midrule
  \multicolumn{2}{r}{\small 接下页} \\
  \endfoot
  \bottomrule
  \endlastfoot
  \texttt{test\_geodesic\_self\_distance}
    & 在 icosphere 上调用 \texttt{geodesic\_distance\_from\_vertex}，
      断言返回 \texttt{Some}；源顶点距离 $< 10^{-6}$；
      至少一个非源顶点距离 $> 0$。 \\
  \texttt{test\_geodesic\_monotonicity}
    & 在 icosphere 上对前若干非源顶点做 sanity check：邻域中至少存在一个顶点
      距离更小，验证距离场单调性。 \\
  \midrule
  \texttt{test\_dijkstra\_self\_distance\_zero}
    & 源点到自身 Dijkstra 距离 $= 0$ \\
  \texttt{test\_dijkstra\_symmetric\_on\_icosphere}
    & icosphere 上 $d(s \to t) \approx d(t \to s)$（容差 $10^{-6}$） \\
  \texttt{test\_dijkstra\_distance\_upper\_bounds\_heat\_method}
    & Dijkstra 与 Heat Method 的最大距离在 3 倍以内（图距离 vs 测地线） \\
  \texttt{test\_dijkstra\_multi\_source}
    & 多源 Dijkstra：每个顶点距离 $\le$ 单源距离；源点距离 $= 0$ \\
  \texttt{test\_dijkstra\_shortest\_path\_self}
    & \texttt{source == target} 时返回 \texttt{[source]} \\
  \texttt{test\_dijkstra\_shortest\_path\_to\_neighbor}
    & 邻居路径长度为 2（\texttt{[source, target]}） \\
  \texttt{test\_dijkstra\_shortest\_path\_consistency\_with\_distance}
    & 路径长度与距离一致：$\text{dist}[\text{target}] \approx \sum \|p_{i+1} - p_i\|$ \\
  \midrule
  \texttt{test\_multi\_source\_geodesic\_sources\_zero}
    & 多源 Heat Method 所有源点距离 $< 10^{-6}$（强制置零生效） \\
  \texttt{test\_multi\_source\_geodesic\_le\_single\_source}
    & 多源最大距离 $\le$ 单源最大距离（容差 $10^{-6}$） \\
  \midrule
  \texttt{mmp\_self\_distance\_zero}
    & MMP 源顶点距离 $= 0$ \\
  \texttt{mmp\_le\_dijkstra}
    & MMP 距离 $\le$ Dijkstra 距离（可穿面，路径更短） \\
  \texttt{mmp\_symmetric}
    & icosphere 上 $d(s \to t) \approx d(t \to s)$（5\% 容差） \\
  \texttt{mmp\_multi\_source}
    & 多源 MMP 距离 $\le$ 单源距离 \\
  \texttt{mmp\_multi\_source\_unified\_matches\_naive}
    & 统一传播与逐源取最小数值一致（容差 $10^{-9}$） \\
  \texttt{mmp\_flat\_grid}
    & 平面网格上 MMP 距离接近欧氏距离（容差 0.1--0.2） \\
  \midrule
  \multicolumn{2}{l}{\textit{边界与无效输入测试（icosphere(0) 与空网格）}} \\
  \texttt{dijkstra\_empty\_mesh\_returns\_empty}
    & 空网格 \texttt{MeshStorage::new()} 上 \texttt{dijkstra\_geodesic} 返回空 \texttt{Vec} \\
  \texttt{mmp\_empty\_mesh\_returns\_empty}
    & 空网格上 \texttt{mmp\_geodesic} 返回空 \texttt{Vec} \\
  \texttt{dijkstra\_multi\_source\_empty\_sources\_returns\_infinity}
    & 非空网格但 \texttt{sources = \&[]} 时返回全 \texttt{INFINITY} \\
  \texttt{dijkstra\_invalid\_source\_returns\_infinity}
    & \texttt{VertexId::default()} 不在网格中，返回全 \texttt{INFINITY} \\
  \texttt{mmp\_invalid\_source\_returns\_infinity}
    & 同上用 \texttt{mmp\_geodesic}，返回全 \texttt{INFINITY} \\
  \texttt{geodesic\_distance\_from\_vertex\_invalid\_source\_returns\_none}
    & 无效 source 经 \texttt{?} 传播返回 \texttt{None} \\
  \texttt{multi\_source\_geodesic\_empty\_sources\_returns\_none}
    & 空源集直接返回 \texttt{None} \\
  \texttt{dijkstra\_with\_parent\_invalid\_source\_returns\_empty}
    & 无效 source 返回 \texttt{(Vec::new(), Vec::new())} \\
  \texttt{dijkstra\_shortest\_path\_invalid\_target\_returns\_empty}
    & 无效 target 返回空 \texttt{Vec} \\
  \texttt{dijkstra\_distance\_nonnegative}
    & 所有 Dijkstra 距离 $\ge 0$ \\
  \texttt{mmp\_multi\_source\_single\_source\_matches\_naive}
    & \texttt{mmp\_multi\_source\_geodesic(\&[v0])} 与 \texttt{mmp\_geodesic(v0)} 数值一致 \\
  \texttt{mmp\_distance\_nonnegative}
    & 所有 MMP 距离 $\ge 0$ \\
  \texttt{shortest\_path\_to\_self\_returns\_self}
    & \texttt{target == source} 时回溯立即终止，路径 $=$ \texttt{[source]} \\
  \texttt{shortest\_path\_invalid\_distance\_returns\_empty}
    & 空距离场 \texttt{\&[]} 时 \texttt{target\_idx >= len} 返回空 \texttt{Vec} \\
\end{longtable}
}

\texttt{src/geodesics.rs} 共 \textbf{31 个}单元测试，全库共 \textbf{567 个}。

\end{document}
